\begin{tikzpicture}[font=\small,line cap=round,line join=round]\begin{scope}[scale=0.4]\draw[otNavy,line width=1pt](0.000,0.000)circle(4.000);\fill[otNavy](0.000,0.000)circle(1.6pt)node[below,geoLabel]{$ A $};\fill[otNavy](6.000,0.000)circle(1.6pt)node[below,geoLabel]{$ B $};\draw[otMuted](0.000,0.000)--(6.000,0.000);\draw[otOrange,dashed,line width=.7pt](6.000,0.000)circle(2.000);\node[above,font=\tiny,inner sep=1.2pt,text=otOrange]at(6.000,2.000){$r=2$};\draw[otGreenDark,dashed,line width=.7pt](6.000,0.000)circle(5.000);\node[above,font=\tiny,inner sep=1.2pt,text=otGreenDark]at(6.000,5.000){$r=5$};\draw[otRed,dashed,line width=.7pt](6.000,0.000)circle(10.000);\node[above,font=\tiny,inner sep=1.2pt,text=otRed]at(6.000,10.000){$r=10$};\end{scope}\end{tikzpicture}\par\vspace{1.5mm}{\otRegular\fontsize{8}{10}\selectfont\color{otMuted}Le cercle de centre A et de rayon 4 cm, et des cercles de centre B de rayons 2, 5 et 10 cm (figure réduite).\par}
